\subsubsection{Base64}
\myindex{Base64}

base64エンコーディングは、バイナリデータをテキスト文字列として転送する必要がある場合に非常に人気があります。

本質的に、このアルゴリズムは3つのバイナリバイトを4つの印刷可能文字にエンコードします：
すべての26個のラテン文字（小文字と大文字の両方）、数字、プラス記号（「+」）、スラッシュ記号（「/」）、
合計64文字です。

base64文字列の特徴的な機能の一つは、しばしば（常にではありませんが）1つまたは2つの\gls{padding}
等号記号（「=」）で終わることです。例えば：

\begin{lstlisting}
AVjbbVSVfcUMu1xvjaMgjNtueRwBbxnyJw8dpGnLW8ZW8aKG3v4Y0icuQT+qEJAp9lAOuWs=
\end{lstlisting}

\begin{lstlisting}
WVjbbVSVfcUMu1xvjaMgjNtueRwBbxnyJw8dpGnLW8ZW8aKG3v4Y0icuQT+qEJAp9lAOuQ==
\end{lstlisting}

等号記号（「=」）はbase64エンコード文字列の途中には決して現れません。

手動エンコーディングの例。
0x00、0x11、0x22、0x33の16進バイトをbase64文字列にエンコードしてみましょう：

\lstinputlisting{digging_into_code/strings/base64_ex.sh}

すべての4バイトをバイナリ形式で配置し、6ビットグループに再グループ化しましょう：

\begin{lstlisting}
|  00  ||  11  ||  22  ||  33  ||      ||      |
00000000000100010010001000110011????????????????
| A  || B  || E  || i  || M  || w  || =  || =  |
\end{lstlisting}

最初の3バイト（0x00、0x11、0x22）は4つのbase64文字（「ABEi」）にエンコードできますが、
最後の1つ（0x33）はエンコードできないため、
2つの文字（「Mw」）を使用してエンコードされ、\gls{padding}記号（「=」）が
最後のグループを4文字にパディングするために2回追加されます。
したがって、すべての正しいbase64文字列の長さは常に4で割り切れます。

\myindex{XML}
\myindex{PGP}
Base64は、バイナリデータをXMLに格納する必要がある場合によく使用されます。
「装甲化された」（つまり、テキスト形式の）PGPキーと署名はbase64を使用してエンコードされます。

一部の人々は文字列を難読化するためにbase64を使用しようとします：
\url{http://blog.sec-consult.com/2016/01/deliberately-hidden-backdoor-account-in.html}
\footnote{\url{http://archive.is/nDCas}}。

\myindex{base64scanner}
任意のバイナリファイルをbase64文字列でスキャンするユーティリティがあります。
そのようなユーティリティの一つはbase64scanner\footnote{\url{https://github.com/DennisYurichev/base64scanner}}です。

\myindex{UseNet}
\myindex{FidoNet}
\myindex{Uuencoding}
\myindex{Phrack}
UseNetとFidoNetではるかに人気があった別のエンコーディングシステムはUuencodingです。
バイナリファイルは今でもPhrack雑誌でUuencode形式でエンコードされています。
ほぼ同じ機能を提供しますが、ファイル名もヘッダーに格納されるという点でbase64とは異なります。

\myindex{Tor}
\myindex{base32}
ちなみに：base64の近親もあります：base32。そのアルファベットには~10桁と~26のラテン文字があります。
その有名な使用例の一つはオニオンアドレス
\footnote{\url{https://trac.torproject.org/projects/tor/wiki/doc/HiddenServiceNames}}です。
例：\url{http://3g2upl4pq6kufc4m.onion/}。
\ac{URL}は大文字小文字が混在するラテン文字を持つことができないため、明らかにTor開発者がbase32を使用した理由です。